\section{Casos de uso del instalador}

Esta sección describe las decisiones que debe tomar el técnico durante una instalación. El objetivo es entregar un sistema operativo sin alterar servicios, dispositivos o credenciales que pertenezcan al cliente.

\subsection{Caso 1: Cliente con Home Assistant existente}
\textbf{Perfil:} \texttt{bridge\_ha}.

Use este escenario cuando el cliente ya opera Home Assistant y desea conservar sus integraciones, automatizaciones y usuarios actuales.
\begin{enumerate}[itemsep=0.15cm]
    \item El técnico selecciona \texttt{bridge\_ha} en el asistente.
    \item HomePilot verifica que Home Assistant responda y solicita una URL accesible junto con un token de larga duración.
    \item El técnico prueba la conexión antes de guardarla.
    \item Los dispositivos aprobados se descubren e importan al inventario propio de HomePilot.
\end{enumerate}

\begin{tcolorbox}[title=Resultado esperado]
HomePilot opera como consola local y bridge. Home Assistant existente no se reemplaza, no se reinicia sin autorización y mantiene sus datos intactos.
\end{tcolorbox}

\subsection{Caso 2: Cliente sin Home Assistant y con necesidad de integraciones}
\textbf{Perfil:} \texttt{ha\_companion}.

Use este escenario cuando el proyecto requiere integraciones compatibles con Home Assistant, pero el cliente aún no dispone de una instalación.
\begin{enumerate}[itemsep=0.15cm]
    \item El técnico selecciona \texttt{ha\_companion}.
    \item El asistente prepara el servicio complementario de Home Assistant junto con HomePilot.
    \item El responsable designado crea la primera cuenta de Home Assistant.
    \item El técnico enlaza HomePilot usando un token de larga duración y valida la conexión.
    \item Se incorporan únicamente las integraciones y dispositivos incluidos en el alcance contratado.
\end{enumerate}

\begin{tcolorbox}[title=Resultado esperado]
El cliente recibe una instalación local completa, con HomePilot como interfaz de uso cotidiano y Home Assistant como capa de integración compatible.
\end{tcolorbox}

\subsection{Caso 3: Instalación nativa de HomePilot}
\textbf{Perfil:} \texttt{native\_only}.

Use este escenario cuando el alcance se limita a funciones nativas disponibles en HomePilot y no se requiere Home Assistant durante la entrega inicial.
\begin{enumerate}[itemsep=0.15cm]
    \item El técnico selecciona \texttt{native\_only}.
    \item El asistente prepara HomePilot sin crear ni enlazar Home Assistant.
    \item El responsable designado crea la primera cuenta administrativa de HomePilot.
    \item Se registra el hogar, las estancias y los dispositivos nativos incluidos en el proyecto.
\end{enumerate}

\begin{tcolorbox}[title=Resultado esperado]
La instalación queda lista para control local. Una integración compatible puede añadirse posteriormente mediante una orden de servicio, sin alterar el perfil original sin confirmación técnica.
\end{tcolorbox}

\subsection{Caso 4: Actualización de una instalación existente}

Use este escenario cuando HomePilot ya está instalado y debe recibir una versión aprobada.
\begin{lstlisting}[language=bash]
git pull --ff-only
bash scripts/homepilot-maintenance.sh --deploy --yes
\end{lstlisting}

El mantenimiento conserva el perfil guardado en \texttt{.env} y reutiliza la base de datos existente. Antes del despliegue, el técnico debe revisar el estado de los servicios y el espacio disponible.

\subsection{Caso 5: Diagnóstico sin cambios}

Use esta opción para revisar conectividad, puertos, espacio y estado de servicios sin modificar archivos ni reiniciar contenedores.
\begin{lstlisting}[language=bash]
bash scripts/install-edge-office.sh --status
\end{lstlisting}

En el asistente guiado, corresponde a la acción \textit{Ejecutar solo diagnóstico}. Es la alternativa correcta antes de intervenir una instalación con incidentes.

\subsection{Caso 6: Integraciones de comunidad}

Use este escenario cuando una integración requerida, como SonoffLAN mediante HACS, no forma parte de Home Assistant base.
\begin{enumerate}[itemsep=0.15cm]
    \item El técnico verifica si HACS ya está instalado.
    \item Solicita autorización antes de instalar software de comunidad.
    \item Realiza la instalación y reinicio requerido por Home Assistant.
    \item Configura la integración con las credenciales del cliente.
    \item Valida el dispositivo y lo importa en HomePilot desde Descubrimiento.
\end{enumerate}

\begin{tcolorbox}[title=Regla de autorización]
HACS, integraciones de comunidad y credenciales de fabricantes se modifican únicamente con autorización del responsable de la instalación. HomePilot no debe instalar ni sustituir estos componentes de forma silenciosa.
\end{tcolorbox}

\subsection{Caso 7: Entrega al cliente}

Al finalizar cualquier perfil, el técnico debe comprobar que la interfaz local abre, que el administrador puede iniciar sesión, que los dispositivos aprobados responden y que el cliente recibe las credenciales iniciales por un canal seguro. El cliente utiliza HomePilot; las tareas de infraestructura y mantenimiento quedan bajo responsabilidad del equipo técnico autorizado.
